\scanpage{454}\mspage{366}
% archive sheet number: 453
pas forcément le cas.

\deleted{Mais $\deleted{K}\unsure{J|K}$ \unsure{étant|est} \unsure{grande|gde} devant $\pi$,}

Posons donc \inserted{$K=K_0$,} $K'=\deleted{\unsure{K}}K_1$, et construisons par induction une suite croissante de sous-catégories $K_i$ ($i\in\unsure{\mathbb{N}|N}$), de $\mathrm{Ob}\,K_i=\mathrm{Ob}\,J$, et $K_i$ \inserted{(pour $i\geq1$)} étant déduite de $K_{i-1}$ en lui rajoutant \deleted{\unsure{des}} une flèche $p_i:e\longrightarrow e$, comme ci-dessus. Ainsi tte double flèche dans un $K_i$ peut s'égaliser dans $K_{i+1}$. Prenons $K_{\infty}^{\deleted{\unsure{\prime|q|,}}}=\unsure{\varinjlim|\underline{\lim}}\,K_i$, $K_{\infty}$ est une sous-catégorie \uline{filtrante} de $\deleted{K}J$, \unsure{contenant} $K$ et de \unsure{cardinal|card} $\leq\pi$, \deleted{\unsure{définissant une co-}} qed. $\big]$

On peut préciser cette construction ainsi\unsure{ ``|.}

\uline{Corollaire} 4.13.39. Sous les conditions de 4.13.38, supposons que $K$ contienne un objet $e$ \unsure{maximal|minimal} dans $\deleted{K}J$. Soit $\Sigma$ l'ens des endomorphismes $p\unsure{/}$de $e$ dans $J$, tels que
pour tte double-flèche \diagram{\begin{tikzcd}[ampersand replacement=\&,column sep=large]
k\arrow[r,shift left=.48ex,"u"]\arrow[r,shift right=.48ex,"v"'] \& e
\end{tikzcd}} \uline{dans} $K$, on
ait $pu=pv$. \deleted{Alors \unsure{a)} il existe \illegible{1},} Alors

\unsure{a)|c)} $\Sigma$ est stable par composition et il existe une partie non vide $\Sigma_{\infty}$ de $\Sigma$, telle que \deleted{pour \unsure{$p$} \illegible{1} \unsure{catégo}} pour \deleted{\unsure{tous}} $p,p'\in\Sigma_0$, existe
